\begin{titlepage}
\color{navy}\sffamily
\noindent\large The Energy Balance Project
\vspace{35pt}
\par\noindent\fontsize{29}{35}\selectfont\bfseries The Energy\par Balance Project
\vspace{20pt}
\par\noindent\fontsize{21}{27}\selectfont Part I: Physical Intuition\par and the EBU Idea
\vspace{20pt}
\par\noindent\large From life-supporting function to an exact\par account of physical action
\vspace{25pt}
\par\noindent\normalsize Revised Gaussian explanatory edition
\vspace{15pt}
\par\noindent\normalfont This book begins with the reason for the project and develops its mathematics patiently. Physical stories, notation, diagrams, worked calculations, guided problems and practice studios form one continuous course. Curiosity and ordinary arithmetic are enough to begin.
\vfill
\noindent\sffamily\Large Konrad Grzyb
\par\normalsize Author and project originator
\vspace{20pt}
\par\small AUTHOR-REVIEW EDITION\par September 2026
\end{titlepage}

\chapter*{Author, rights, and scientific scope}
Konrad Grzyb is the author and originator of the Energy Balance Project. This revised manuscript was prepared with AI assistance under his direction. It is supplied for author review; completion of the typesetting is not a claim of author acceptance, independent peer review or publication.

The software repository has a separate MIT licence. That software licence does not automatically grant rights in the books. This revision introduces no new rights grant. The original unified explanatory edition is preserved, as are Parts II and III. Their historical findings remain findings about their declared models.

\begin{lawbox}{Scientific boundary}
Definitions, conditional mathematical consequences, implementation checks and observed scientific results are different kinds of statement. This book explains a declared Gaussian potential and derives exact identities. It does not report an experiment establishing recovery, fairness, planetary homeostasis or economic adoption.
\end{lawbox}

The original Part I is the editorial baseline. Its original editable manuscript was unavailable in the inspected project files, so this is a reconstructed, editable revision. The companion revision ledger identifies retained subjects, rewritten examples and removed model-specific material. The shorter introductory candidate remains a separate preserved draft; it is not this book's scale or layout template.

All numbered tank calculations are illustrative mathematics on explicitly supplied states. They are not measurements of hospitals or reservoirs, recommended operating quantities, registered experimental parameters, or model trajectories. A reference or scale in such an example is a teaching choice. A deployed system would require measurement, calibration, institutions and separate evidence.

\chapter*{Preface: the question behind the project}
A clinic can have an orderly financial account while its water system is failing. The reverse is also possible: its equipment works, but a patient cannot afford access. These are distinct failures, and neither disappears because the other is measured well. The motivation of this book is to bring physical function into sharper view without losing that distinction.

Money coordinates claims, obligations and exchanges. It can support investment, maintenance, cooperation and public provision. A monetary record does not, by itself, reconstruct the water withdrawn, electricity transformed, equipment worn, resources displaced or service delivered in an action. An additional physical account might help people ask better questions about these effects. Whether a particular account actually improves decisions is a question for investigation.

The Energy Balance Project proposes a way to value a specified physical change relative to a declared representation of functional balance. It begins with a state: what is present, where it is and in which units. It declares a reference, a scale and a potential measuring departure from that reference. Then it compares the same represented world before and after a finite action. The comparison has a sign. Reducing the declared burden produces positive field value; increasing it produces negative field value.

There is no moral score for the person in this equation. A medicine delivery may be essential and physically burdensome. A perfectly calculated tank redistribution may fail to represent water quality. A useful repair may create a benefit the current state description does not yet contain. Throughout the book, mathematical exactness is paired with a question about representation: exact relative to what?

This revision retains the learning method of the original Part I. We draw a physical scene before writing its matrix. We read subscripts and transposes aloud before using them. We calculate the same action by more than one route. We include zero and negative cases. The studios at the end repeat the entire chain through longer worked problems rather than assuming that recognizing an equation is the same as understanding it.

The central change is the introductory potential. Earlier editions taught one threshold-based model in detail. Here the main example uses a fixed reference and a positive scale to build a quadratic, Gaussian-shaped deviation geometry. Biological threshold laws, lower and upper homeostatic bands, and a particular preservation controller are no longer prerequisites for understanding EBU. Their original theorems remain scoped to the historical models in which they were established.

This change does not make all living systems Gaussian. It gives us a transparent, differentiable example whose finite differences can be calculated exactly. It also makes an important limitation visible: a mathematical landscape does not contain its own actor. Someone or something must choose and perform an action. A field that measures change is not automatically a controller that restores balance.

The book can therefore be read before new scientific tests. Definitions and conditional proofs can be explained now. Recovery, persistence, useful incentives and real institutional effects require their own evidence. The final chapter returns to those questions after the reader has the tools to state them precisely.

\chapter*{The route through this volume}
The opening seven chapters develop the reason for the project: monetary signals, scarcity, unequal exposure, planetary pressures, homeostasis and equilibrium. They distinguish a motivating concern from a conclusion about EBU. None asks the reader to accept a replacement economy in advance.

Chapters 8--16 build the physical and mathematical objects. We follow a pump, map cells and edges, check conservation, declare references and scales, construct the potential, interpret its slope, and calculate the full effect of a finite action. Permission is explained before the exact quote is used, because a favourable value cannot make an impossible movement possible.

Chapters 17--23 explain records and composition. They separate a prediction from verified occurrence, a local calculation from a global audit, an exact group total from an individual attribution, and a route from an instantaneous remote effect. These chapters preserve the original volume's concern with real needs, repairs, privacy, accountability and institutional choice.

Chapters 24--31 are practice studios. Keep a notebook open. Write the units next to every value, calculate before consulting an answer and explain each sign in words. A recurring Ridge--Vale example provides continuity; variations show which conclusions depend on its special numbers. Chapter 32 closes the course with the mathematical statements, their assumptions and the research questions still open.

The combined Part II continues directly at Chapter 33, joining mathematical derivation, implementation and evidence into one course. It develops group value before attribution, attribution before capacity accounts, and physical permission before execution. Detailed historical proofs and experiments sit within the relevant chapters, with their assumptions clearly marked. References to \emph{historical Part II} or \emph{historical Part III} still identify the preserved original editions, not the new numbering. Their threshold models must not be substituted into Gaussian examples without changing the declaration.

The forward series now has eight parts: I introduces the idea; II joins mathematics, implementation and evidence; III develops measurement and falsification; IV studies homeostasis and stochastic dynamics; V develops multiple-action structure; VI addresses distance; VII addresses coordination and memory; VIII addresses institutions. The later six are a research and teaching plan, not established findings.

\chapter*{How to read an equation without fear}
Begin with the nouns. In a statement such as $x_i=18\X$, $x_i$ is the quantity at one named location, $18$ is its numerical value and $\X$ is its stock unit. The subscript identifies a coordinate. It is not a multiplication by $i$. A reference $x_i^*$ is a declared comparison value; the star is a label, not an instruction to multiply.

Read a subtraction as a comparison. If before-burden is 16 and after-burden is 4, then $16-4=12$ means the represented burden decreased by twelve units of the declared account. The arithmetic alone does not identify a payer, establish a wage, or certify that all human needs were satisfied.

A vector is a list whose order matters. A matrix is a rule for combining entries in lists. A derivative describes how a function changes nearby. An integral adds those changing nearby contributions along a specified path. These descriptions will be made precise through small examples before the general forms are needed.

When an equation has a condition beside it, read the condition as part of the equation. A claim for fixed references and scales is not automatically true if either is revised halfway through an account. A proof for complete state records cannot detect damage that the state never represented. A proof for a common action path does not identify a physical path that was never observed.

\begin{intuitionbox}{Four questions beside every calculation}
What physical object does each symbol name? Which quantities are held fixed? What operation is being performed? What conclusion follows, and which conclusion would need additional evidence?
\end{intuitionbox}

Use the glossary when a familiar letter changes role. In particular, stock, scale, potential, marginal, action quantity, attribution and spending capacity are different objects. Similar typography is never permission to exchange their meanings. The same discipline that helps a first-time reader also protects a scientific calculation.
